% Two complementary geometric properties of a conservative force.
% Native size: 11.52 cm wide, matching the divergence/curl figure.
\begin{tikzpicture}[
  x=1cm,
  y=1cm,
  font=\footnotesize,
  line cap=round,
  line join=round,
  panel title/.style={font=\small,inner sep=0pt},
  label/.style={inner sep=0pt},
  vector arrow/.style={
    draw=black,
    line width=0.5pt,
    -{Latex[length=1.4mm,width=1.05mm]}
  },
  path curve/.style={
    vector arrow,
    shorten <=1.6pt,
    shorten >=1.6pt
  },
  equipotential/.style={draw=black,line width=0.5pt}
]
  \path[use as bounding box] (0,0) rectangle (11.52,4.25);

  % (a) Both paths run from A to B; their work integrals agree.
  \begin{scope}[shift={(2.60,2.20)}]
    \node[panel title] at (0,1.75) {(a) Path independence};
    \coordinate (A) at (-1.90,0);
    \coordinate (B) at (1.90,0);
    \draw[path curve]
      (A) .. controls (-1.02,1.08) and (1.02,1.08) .. (B);
    \draw[path curve]
      (A) .. controls (-1.02,-1.08) and (1.02,-1.08) .. (B);
    \fill[black] (A) circle[radius=1.2pt] node[left=4pt] {$A$};
    \fill[black] (B) circle[radius=1.2pt] node[right=4pt] {$B$};
    \node[label] at (0,1.09) {$\mathcal C_1$};
    \node[label] at (0,-1.09) {$\mathcal C_2$};
    \node[label] at (0,-1.86) {$\displaystyle
      \int_{\mathcal C_1}\!\mathbf F\!\cdot\!\mathrm{d}\mathbf r
      =
      \int_{\mathcal C_2}\!\mathbf F\!\cdot\!\mathrm{d}\mathbf r$};
  \end{scope}

  % (b) U increases upwards; the force opposes its normal gradient.
  \begin{scope}[shift={(8.92,2.20)}]
    \node[panel title] at (0,1.75) {(b) Potential field};
    \foreach \cfY in {-0.68,-0.34,0,0.34,0.68,1.02}{
      \draw[equipotential] (-1.90,\cfY) -- (1.90,\cfY);
    }
    \node[label,anchor=east] at (-2.05,-0.68) {$U_1$};
    \node[label,anchor=east] at (-2.05,0) {$U_2$};
    \node[label,anchor=east] at (-2.05,1.02) {$U_3$};

    \coordinate (P) at (0,0);
    \draw[vector arrow] (P) -- (0,1.24);
    \draw[vector arrow] (P) -- (0,-1.24);
    % Labels sit beyond the equipotentials, so no line crosses the text.
    \node[label,anchor=west] at (0.18,1.24) {$\nabla U$};
    \node[label,anchor=west] at (0.18,-1.24) {$\mathbf F$};
    \draw[black,line width=0.5pt]
      (0.14,0) -- (0.14,0.14) -- (0,0.14);
    \fill[black] (P) circle[radius=1.2pt];
    \node[label] at (0,-1.64) {$U_1<U_2<U_3$};
    \node[label] at (0,-2.02) {$\mathbf F=-\nabla U$};
  \end{scope}
\end{tikzpicture}
